% ============================================================
% section4_database.tex - Раздел 4: Пакетная загрузка и хранение данных
% ============================================================
% !TeX program = xelatex
% Компиляция: xelatex --shell-escape main.tex

\section{Система хранения данных: Пакетная загрузка и СУБД}\label{sec:database}

В предыдущем разделе мы успешно загружали и строили графики для \textit{одного} сигнала и одного фона. Но реальный физический эксперимент выглядит иначе: спектрометр последовательно поворачивает поляризаторы и записывает отдельный файл для каждого угла, делая при этом по несколько повторных проходов (\textit{реплик}), чтобы минимизировать случайные погрешности.

Прописывать пути ко всем файлам вручную означает верный путь к хаосу: раздутый код, опечатки в путях, потеря связи между парами «сигнал-фон». Наша цель в этом разделе --- написать модуль, который автоматически просканирует папку с результатами, считает все данные и сохранит их в единую структурированную базу данных прямо в оперативной памяти.

\subsection{Архитектура СУБД в оперативной памяти}\label{ssec:db_arch}

В типичных студенческих скриптах обработка данных часто превращается в хаос: создаются десятки несвязанных массивов --- \mycode{signals_raw}, \mycode{backgrounds_raw}, \mycode{transmissions}. При этом связь между файлом сигнала образца и соответствующим ему фоном теряется или жестко зашивается через индексы, что приводит к критическим ошибкам при добавлении новых измерений.

Чтобы решить эту проблему раз и навсегда, мы применим профессиональный архитектурный паттерн --- \textbf{систему управления базами данных (СУБД)}. В нашем случае это будет простая и быстрая реляционная (связанная) база данных, живущая прямо в оперативной памяти компьютера в виде единого централизованного словаря \mycode{data_store} (хэш-таблицы).

Словарь в Python хранит данные в виде пар <<ключ --- значение>>. В качестве \textit{значения} мы записываем кортеж двух одномерных массивов NumPy \mycode{(t, E)}, а в качестве \textit{ключа} --- структурированные метаданные измерения.

\subsubsection{Схема составных ключей}

Ключом в нашем словаре-хранилище \mycode{data_store} будет строго типизированный кортеж (tuple --- неизменяемый упорядоченный список) из четырех элементов, однозначно идентифицирующий любое измерение:
\begin{equation}
\text{Key} = (\text{Tag}, \text{Angle}_1, \text{Angle}_2, \text{Repetition})
\end{equation}

Каждый элемент имеет чёткий физический и логический смысл:
\begin{enumerate}
    \item \textbf{\mycode{Tag} (строка):} Уникальный идентификатор типа данных:
    \begin{itemize}
        \item \mycode{'signal_raw'} --- сырой временной ТГц-сигнал образца для конкретного прохода.
        \item \mycode{'bg_raw'} --- сырой временной фоновый (опорный) сигнал для конкретного прохода.
        \item \mycode{'spec_avg'} --- итоговый усредненный спектр пропускания по мощности (после обработки всех реплик).
        \item \mycode{'spec_bg_avg'} --- итоговый усредненный спектр мощности опорного сигнала.
    \end{itemize}
    \item \textbf{\mycode{Angle1} (число с плавающей точкой):} Нормализованный угол установки первого ротатора (поляризатора) в градусах.
    \item \textbf{\mycode{Angle2} (число с плавающей точкой):} Угол установки второго ротатора (анализатора) в градусах.
    \item \textbf{\mycode{Repetition} (целое число или \mycode{None}):} Номер повторного измерения. Для усредненных данных, где реплики уже объединены, этот параметр равен \mycode{None}.
\end{enumerate}

Пример ключа для первого прохода сигнала образца при угле поляризатора $10^\circ$: \mycode{('signal_raw', 10.0, 0.0, 1)}. Такая структура гарантирует, что мы никогда не перепутаем сигнал с фоном, разные углы или разные повторения.

\subsubsection{Декларативная SQL-подобная фильтрация}

Одно из главных преимуществ такого СУБД-подхода --- невероятно простая фильтрация данных с помощью встроенных в Python генераторов словарей (dict comprehensions --- компактный синтаксис сборки новых коллекций на лету). Например, извлечь из базы все усредненные спектры образцов для построения графиков можно всего одной строчкой:
\begin{minted}[linenos=false]{python}
spectra_to_plot = {
    k: v
    for k, v in data_store.items()
    if k[0] == 'spec_avg'
}
\end{minted}

Конструкция \mycode{data_store.items()} возвращает пары <<ключ --- значение>>. Наш фильтр проверяет первый элемент ключа \mycode{k[0]} (строку-тег) и оставляет только записи с тегом \mycode{'spec_avg'}.

\subsection{Разработка загрузчика базы данных}\label{ssec:load_dataset}

Условия проведения физического эксперимента непрерывно меняются во времени: колеблется температура и влажность воздуха, дрейфует средняя мощность накачки лазера. По этой причине каждый замер сигнала образца жестко привязан к конкретному файлу фона, измеренному \textit{в той же серии проходов}.

\warning{\textbf{Критическая физическая ошибка:} Абсолютно недопустимо сначала усреднить отдельно все сигналы образца и отдельно все фоны во временной области, а затем делить их спектры друг на друга. Единственный физически корректный алгоритм:\\1. Рассчитать спектры сигнала и парного ему фона строго в рамках одного прохода.\\2. Разделить их мощности: $T_i(\nu) = |E_{\text{sig},i}(\nu)|^2 / |E_{\text{bg},i}(\nu)|^2$. Деление мгновенно скомпенсирует дрейф мощности лазера!\\3. И только затем усреднить пропускания по повторениям: $T_{\text{avg}}(\nu) = \frac{1}{N}\sum T_i(\nu)$.\\Наша СУБД-структура обязана строго поддерживать этот алгоритм.}

Добавим функцию \mycode{load_dataset_to_store} в наш модуль \mycode{data_loader.py}. Полный листинг этого модуля для сверки доступен в \textbf{\textit{Приложении~\ref{app:data_loader}}}.

\begin{listing}[H]
	\caption{Реляционная загрузка данных в базу (\mycode{data_loader.py})}\label{lst:load_dataset}
	\begin{minted}{python}
def load_dataset_to_store(data_dir) -> dict:
    data_store = {}
    path = Path(data_dir)

    # Шаг 1. Предварительно загружаем все фоновые файлы
    bg_files = {}
    for filepath in path.glob("*.txt"):
        if filepath.name.startswith("bg"):
            bg_id = filepath.stem   # например, 'bg_1'
            t, E = load_tds_data(filepath)
            bg_files[bg_id] = (t, E)

    # Шаг 2. Сканируем сигналы образцов и связываем их с фонами
    for filepath in path.glob("*.txt"):
        filename = filepath.name
        if filename.startswith("bg"):
            continue

        try:
            angle1, angle2, rep, bg_name = parse_filename(filename)
        except ValueError:
            continue

        t, E_signal = load_tds_data(filepath)

        # Сохраняем сырой сигнал под составным ключом
        data_store[('signal_raw', angle1, angle2, rep)] = (t, E_signal)

        # Привязываем соответствующий фоновый сигнал
        if bg_name in bg_files:
            t_bg, E_bg = bg_files[bg_name]
            data_store[('bg_raw', angle1, angle2, rep)] = (t_bg, E_bg)

    return data_store
	\end{minted}
\end{listing}

\subsubsection{Разбор кода: Пакетный загрузчик}

\begin{syntaxdesc}
    \item[Инициализация базы данных \mycode{data_store = \{\}} (строка 2)] Создаёт пустой словарь-хранилище, в котором будут жить все экспериментальные массивы на время работы программы. Такой подход называют <<базой данных в оперативной памяти>> (in-memory database).
    \item[Сканирование директории \mycode{path.glob("*.txt")} (строки 7, 14)] Метод объекта \mycode{Path} возвращает итератор по всем файлам с расширением \mycode{.txt} в указанной папке, избавляя нас от ручного перечисления.
    \item[Критерий фонового файла \mycode{filepath.name.startswith("bg")} (строки 8, 16)] Проверяет, начинается ли имя файла с букв \mycode{"bg"}. На Шаге~1 это позволяет собрать все фоновые файлы в промежуточный словарь \mycode{bg_files}, а на Шаге~2 --- пропустить их, чтобы обрабатывать только образцы.
    \item[Получение имени без расширения \mycode{filepath.stem} (строка 9)] Атрибут \mycode{stem} объекта \mycode{Path} возвращает имя файла без суффикса: например, для \mycode{'bg_1.txt'} вернёт \mycode{'bg_1'}. Это имя мы используем как ключ в промежуточном словаре \mycode{bg_files}.
    \item[Пропуск нераспознанных файлов \mycode{except ValueError: continue} (строки 21--22)] Если \mycode{parse_filename} не может разобрать имя файла (например, это служебный или системный файл), функция бросает исключение \mycode{ValueError}. Мы его перехватываем и просто переходим к следующему файлу, не прерывая весь цикл.
    \item[Запись в базу \mycode{data_store[('signal_raw', angle1, angle2, rep)] = (t, E_signal)} (строки 27, 32)] Обрати внимание: и для сигнала, и для фона записываются \textit{одинаковые} углы и номер повторения. Это создаёт неразрывную связь между парой и позволяет впоследствии мгновенно найти фон по ключу сигнала, просто заменив тег с \mycode{'signal_raw'} на \mycode{'bg_raw'}.
\end{syntaxdesc}

\note{Программа сканирует папку в два прохода. На первом шаге она загружает все референсные сигналы в словарь \mycode{bg_files}. На втором --- разбирает имена файлов образца, извлекает углы и номер прохода, а затем формирует в \mycode{data_store} неразрывные пары с идентичными метаданными в ключах.}

\subsection{Тестирование пакетной загрузки и выборка данных}\label{ssec:db_test}

Теперь, когда вся папка данных автоматически сканируется и структурируется, давай напишем проверочный скрипт. В Разделе~\ref{sec:loading} мы загружали конкретные файлы сигналов образца и фона вручную, жестко прописывая их полные пути. Теперь все файлы лежат в нашей базе данных \mycode{data_store} в виде кортежей по составным ключам. 

Чтобы извлекать нужную пару (сигнал образца и соответствующий ему фон) красиво и без прямого ручного обращения к словарю по сложным ключам, мы напишем вспомогательную функцию-запрос \mycode{get_signal_pair} в модуле \mycode{utils.py}. Она будет выполнять роль простого фильтра (подобие SQL-запроса) для сборки необходимых векторов. Добавь эту функцию в свой файл \mycode{utils.py}:

\begin{listing}[H]
	\caption{Функция извлечения пары сигналов из базы данных (\mycode{utils.py})}\label{lst:get_signal_pair}
	\begin{minted}{python}
def get_signal_pair(db: dict, angle1: float, angle2: float, rep: int) -> tuple[np.ndarray, np.ndarray, np.ndarray, np.ndarray]:
    """
    Извлекает из базы данных db пару сигналов (образца и фона) по углам и реплике.
    Возвращает кортеж (t_bg, E_bg, t_sig, E_sig).
    """
    sig_key = ('signal_raw', angle1, angle2, rep)
    bg_key = ('bg_raw', angle1, angle2, rep)
    
    if sig_key not in db:
        raise KeyError(f"Ключ сигнала {sig_key} отсутствует в базе данных.")
    if bg_key not in db:
        raise KeyError(f"Ключ фона {bg_key} отсутствует в базе данных.")
        
    t_sig, E_sig = db[sig_key]
    t_bg, E_bg = db[bg_key]
    
    return t_bg, E_bg, t_sig, E_sig
	\end{minted}
\end{listing}

\subsubsection[Разбор кода: Функция-запрос get\_signal\_pair]{Разбор кода: Функция-запрос \mycode{get_signal_pair}}
\begin{syntaxdesc}
    \item[Составные ключи \mycode{sig_key} и \mycode{bg_key} (строки 6--7)] Динамически формируют кортежи-ключи для поиска в базе данных по переданным параметрам углов и номеру реплики.
    \item[Защитная проверка \mycode{if sig_key not in db} (строки 9--12)] Проверяет наличие ключей в базе данных и выбрасывает понятное исключение \mycode{KeyError}, если файлы по ошибке были удалены или не загрузились.
    \item[Извлечение и возврат данных \mycode{return t_bg, E_bg, t_sig, E_sig} (строки 14--17)] Считывает кортежи временной сетки и амплитуды и возвращает единый плоский кортеж из четырех массивов для передачи в модуль визуализации.
\end{syntaxdesc}

Теперь обновим главный управляющий файл (скрипт-оркестратор) \mycode{main.py}. Вместо сырой ручной загрузки файлов из Раздела~\ref{sec:loading} мы проинициализируем базу данных, сделаем SQL-подобный запрос для той же пары углов ($50^\circ$ и $0^\circ$, реплика $1$) и передадим извлеченные векторы в готовую функцию визуализации:

\begin{listing}[H]
	\caption{Тест инициализации базы данных и визуализации (\mycode{main.py})}\label{lst:test_db_load}
	\begin{minted}{python}
import config
from data_loader import load_dataset_to_store
from utils import get_signal_pair
from plots import plot_time_signals

def main():
    # 1. Инициализация базы данных в оперативной памяти
    db = load_dataset_to_store(config.DATA_DIR)
    
    # 2. SQL-подобный запрос для извлечения пары сигналов
    t_bg, E_bg, t_sig, E_sig = get_signal_pair(db, angle1=50.0, angle2=0.0, rep=1)
    
    # 3. Визуализация извлеченной пары во временной области
    plot_time_signals(t_bg, E_bg, t_sig, E_sig)

if __name__ == "__main__":
    main()
	\end{minted}
\end{listing}

Запусти этот скрипт в PyCharm. На экране должно появиться точно такое же интерактивное окно визуализации временных ТГц-сигналов, как и на Рисунке~\ref{fig:3_1}, но теперь эти данные загружены из нашей базы данных!

Попробуй поэкспериментировать: поменяй значения углов или номер реплики в вызове \mycode{get_signal_pair} (например, на \mycode{angle1=10.0, angle2=0.0, rep=2}) и посмотри, как изменятся графики временных сигналов. Это демонстрирует гибкость нашей структуры: данные отделены от логики отрисовки.

\begin{minitaskbox}
	\textbf{Мини-задание: Вывод списка доступных углов} \\
	Измени скрипт \mycode{main.py} так, чтобы он не просто выводил суммарное количество файлов, а формировал список всех уникальных углов первого поляризатора, найденных в загруженной базе данных. \\
	\\
	\textbf{Подсказка:} Пройдись по ключам базы, отфильтруй по тегу \mycode{'signal_raw'}, извлеки второй элемент каждого ключа (угол) и сохрани в множество \mycode{set()}, чтобы убрать дубликаты повторных проходов.
\end{minitaskbox}

\subsection{Итоги раздела}

В этом разделе мы заложили фундамент для работы со всем массивом экспериментальных данных:
\begin{enumerate}
    \item Спроектировали локальную реляционную СУБД в оперативной памяти на основе словаря с составными ключами \mycode{(Tag, Angle1, Angle2, Repetition)}.
    \item Написали пакетный загрузчик \mycode{load_dataset_to_store}, автоматически сканирующий папку и связывающий сигналы образца с соответствующими фонами.
    \item Убедились, что архитектура базы данных строго поддерживает физически корректный алгоритм попарного вычитания дрейфа прибора.
\end{enumerate}

Теперь, когда все данные загружены и структурированы, в следующем разделе мы применим цифровую обработку сигналов к каждой паре и в пакетном режиме рассчитаем спектры пропускания для всех углов!
